% from ega1-6-fr.tex
\oldpage[I]{145}
On peut donc dire que la notion de préschéma de type fini sur $Y$ est
\emph{locale sur $Y$}.

\begin{proposition}[6.3.3]
\label{I.6.3.3-fr}
Soient $X$, $Y$ deux schémas affines~; pour que $X$ soit de type fini sur
$Y$, il faut et il suffit que $A(X)$ soit une algèbre de type fini sur
$A(Y)$.
\end{proposition}

La condition étant évidemment suffisante, prouvons qu'elle est nécessaire.
Posons $A=A(Y)$, $B=A(X)$~; en vertu de
(\hyperref[I.6.3.2-fr]{6.3.2}), il existe un recouvrement ouvert affine fini
$(V_i)$ de $X$ tel que chacun des anneaux $A(V_i)$ soit une $A$-algèbre de
type fini. En outre, les $V_i$ étant quasi-compacts, on peut recouvrir chacun
d'eux par un nombre fini d'ouverts de la forme $D(g_{ij})\subset V_i$, où
$g_{ij}\in B$~; si $\varphi_i$ est l'homomorphisme $B\to A(V_i)$ correspondant
à l'injection canonique $V_i\to X$, on a
\[
  B_{g_{ij}}=(A(V_i))_{\varphi_i(g_{ij})}
  =A(V_i)[1/\varphi_i(g_{ij})],
\]
donc $B_{g_{ij}}$ est une $A$-algèbre de type fini. On peut donc se ramener
au cas où $V_i=D(g_i)$ avec $g_i\in B$. Par hypothèse, il existe une partie
finie $F_i$ de $B$ et un entier $n_i\geq 0$ tels que $B_{g_i}$ soit l'algèbre
engendrée sur $A$ par les éléments $b_i/g_i^{n_i}$, où $b_i$ parcourt $F_i$.
Comme les $g_i$ sont en nombre fini, on peut d'ailleurs supposer tous les
$n_i$ égaux à un même entier $n$. En outre, comme les $D(g_i)$ forment un
recouvrement de $X$, l'idéal engendré dans $B$ par les $g_i$ est égal à $B$,
autrement dit, il existe des $h_i\in B$ tels que
$\sum_i h_i g_i=1$. Soit alors $F$ la partie finie de $B$, réunion des $F_i$,
de l'ensemble des $g_i$ et de l'ensemble des $h_i$~; montrons que le
sous-anneau $B'=A[F]$ de $B$ est égal à $B$. Par hypothèse, pour tout $b\in B$
et tout $i$, l'image canonique de $b$ dans $B_{g_i}$ est de la forme
$b'_i/g_i^{m_i}$, où $b'_i\in B'$~; en multipliant les $b'_i$ par des
puissances convenables des $g_i$, on peut encore supposer tous les $m_i$
égaux à un même entier $m$. Par définition des anneaux de fractions, il y a
donc un entier $N$ (dépendant de $b$) tel que $N\geq m$ et $g_i^N b\in B'$ pour
tout $i$~; or, \emph{dans l'anneau $B'$}, les $g_i^N$ engendrent l'idéal $B'$,
puisqu'il en est ainsi des $g_i$ (les $h_i$ appartenant à $B'$)~; il y a donc
des $c_i\in B'$ tels que $\sum_i c_i g_i^N=1$, d'où
$b=\sum_i c_i g_i^N b\in B'$, C.Q.F.D.

\begin{proposition}[6.3.4]
\label{I.6.3.4-fr}
\begin{enumerate}
  \item[(i)] Toute immersion fermée est de type fini.
  \item[(ii)] Le composé de deux morphismes de type fini est de type fini.
  \item[(iii)] Si $f:X\to X'$, $g:Y\to Y'$ sont deux $S$-morphismes de type
    fini, $f\times_S g$ est de type fini.
  \item[(iv)] Si $f:X\to Y$ est un $S$-morphisme de type fini, $f_{(S')}$ est
    de type fini pour toute extension $g:S'\to S$ du préschéma de base.
  \item[(v)] Si le composé $g\circ f$ de deux morphismes est de type fini, et
    si $g$ est séparé, $f$ est de type fini.
  \item[(vi)] Si un morphisme $f$ est de type fini, il en est de même de
    $f_{\mathrm{red}}$.
\end{enumerate}
\end{proposition}

En vertu de (\hyperref[I.5.5.12-fr]{5.5.12}), il suffit de démontrer (i),
(ii) et (iv).

Pour établir (i), on peut se borner au cas d'une injection canonique
$X\to Y$, $X$ étant un sous-préschéma fermé de $Y$~; en outre
(\hyperref[I.6.3.2-fr]{6.3.2}), on peut supposer $Y$ affine, auquel cas $X$
est aussi affine (\hyperref[I.4.2.3-fr]{4.2.3}) et son anneau est isomorphe à
un anneau quotient $A/\mathfrak{J}$, où $A$ est l'anneau de $Y$ et
$\mathfrak{J}$ un idéal de $A$~; comme $A/\mathfrak{J}$ est de type fini sur
$A$, la conclusion en résulte.

Démontrons maintenant (ii). Soient $f:X\to Y$, $g:Y\to Z$ deux morphismes de
type fini, et soit $U$ un ouvert affine de $Z$~; $g^{-1}(U)$ admet un
recouvrement fini par des ouverts affines $V_i$ tels que $A(V_i)$ soit une
algèbre de type fini sur $A(U)$ (\hyperref[I.6.3.2-fr]{6.3.2})~; de même,

